% chuong_v3/fig_ch11_viterbi_trellis.tex
% Sơ đồ đồ hình Trellis của giải thuật Viterbi trong Praat Path Finder
% Đối chiếu giữa quy hoạch động toàn cục (Praat) và bộ lọc cục bộ Median-3 (FPGA)
\resizebox{\linewidth}{!}{%
\begin{tikzpicture}[
    >=Stealth,
    scale=0.88,
    frame_box/.style={rectangle, draw=blue!70!black, fill=blue!3, rounded corners=3pt, minimum width=1.6cm, minimum height=5.2cm, inner sep=2pt},
    cand_uv/.style={circle, draw=red!80!black, fill=red!10, thick, minimum size=0.68cm, inner sep=0pt, font=\scriptsize\bfseries},
    cand_v/.style={circle, draw=blue!80!black, fill=blue!10, thick, minimum size=0.68cm, inner sep=0pt, font=\scriptsize\bfseries},
    cand_best/.style={circle, draw=orange!90!black, fill=orange!25, very thick, minimum size=0.72cm, inner sep=0pt, font=\scriptsize\bfseries},
    wire_trans/.style={draw=black!30, thin, ->},
    wire_viterbi/.style={draw=red!85!black, very thick, ->},
    lbl_tag/.style={font=\scriptsize, fill=white, inner sep=2pt, rounded corners=2pt}
]

    % -------------------------------------------------------------------------
    % CÁC KHUNG THỜI GIAN TRONG ĐỒ HÌNH TRELLIS
    % -------------------------------------------------------------------------
    % Khung t-1
    \node[frame_box] (f_prev) at (2.0, 2.5) {};
    \node[below=7pt, font=\scriptsize\bfseries, text=blue!80!black] at (f_prev.south) {Khung $t-1$};

    % Khung t
    \node[frame_box] (f_cur) at (6.0, 2.5) {};
    \node[below=7pt, font=\scriptsize\bfseries, text=blue!80!black] at (f_cur.south) {Khung $t$};

    % Khung t+1
    \node[frame_box] (f_next) at (10.0, 2.5) {};
    \node[below=7pt, font=\scriptsize\bfseries, text=blue!80!black] at (f_next.south) {Khung $t+1$};

    % Khung N (Khung cuối cùng)
    \node[frame_box] (f_end) at (14.0, 2.5) {};
    \node[below=7pt, font=\scriptsize\bfseries, text=blue!80!black] at (f_end.south) {Khung cuối $N$};

    % Dấu chấm lửng giữa t+1 và N
    \node[font=\large\bfseries, text=black!60] at (12.0, 2.5) {$\cdots$};

    % -------------------------------------------------------------------------
    % CÁC ỨNG VIÊN TRONG TỪNG KHUNG (15 ứng viên mỗi khung)
    % -------------------------------------------------------------------------
    % Khung t-1
    \node[cand_uv] (c_p_uv) at (2.0, 4.4) {$c_1$};
    \node[cand_best] (c_p_v1) at (2.0, 3.1) {$c_2$};
    \node[cand_v] (c_p_v2) at (2.0, 1.8) {$c_3$};
    \node[font=\scriptsize, text=black!60] at (2.0, 0.7) {$\vdots$};
    \node[cand_v] (c_p_vm) at (2.0, 0.1) {$c_{15}$};

    % Khung t
    \node[cand_uv] (c_c_uv) at (6.0, 4.4) {$c_1$};
    \node[cand_best] (c_c_v1) at (6.0, 3.1) {$c_2$};
    \node[cand_v] (c_c_v2) at (6.0, 1.8) {$c_3$};
    \node[font=\scriptsize, text=black!60] at (6.0, 0.7) {$\vdots$};
    \node[cand_v] (c_c_vm) at (6.0, 0.1) {$c_{15}$};

    % Khung t+1
    \node[cand_uv] (c_n_uv) at (10.0, 4.4) {$c_1$};
    \node[cand_best] (c_n_v1) at (10.0, 3.1) {$c_2$};
    \node[cand_v] (c_n_v2) at (10.0, 1.8) {$c_3$};
    \node[font=\scriptsize, text=black!60] at (10.0, 0.7) {$\vdots$};
    \node[cand_v] (c_n_vm) at (10.0, 0.1) {$c_{15}$};

    % Khung N
    \node[cand_best] (c_e_uv) at (14.0, 4.4) {$c_1$};
    \node[cand_v] (c_e_v1) at (14.0, 3.1) {$c_2$};
    \node[cand_v] (c_e_v2) at (14.0, 1.8) {$c_3$};
    \node[font=\scriptsize, text=black!60] at (14.0, 0.7) {$\vdots$};
    \node[cand_v] (c_e_vm) at (14.0, 0.1) {$c_{15}$};

    % -------------------------------------------------------------------------
    % CÁC ĐƯỜNG CHUYỂN TRẠNG THÁI VÀ CHI PHÍ (TRANSITION COSTS)
    % -------------------------------------------------------------------------
    % 1. Chuyển tiếp vô thanh - vô thanh: UV -> UV (chi phí = 0)
    \draw[wire_trans, draw=green!60!black] (c_p_uv.east) -- (c_c_uv.west);
    \node[font=\scriptsize, fill=white, inner sep=2pt, rounded corners=2pt, draw=green!60!black, text=green!60!black] at (4.0, 4.95) {UV--UV: $0$};

    % 2. Đổi trạng thái V/UV: c_1 (vô thanh) -> c_2 (hữu thanh) (chi phí = 0,14)
    \draw[wire_trans, draw=purple!80!black] (c_p_uv.east) -- (c_c_v1.west);
    \path (c_p_uv.east) -- (c_c_v1.west) node[pos=0.48, sloped, above=3pt, font=\scriptsize, fill=white, inner sep=2pt, rounded corners=2pt, draw=purple!60, text=purple!80!black] {Đổi V/UV: $0{,}14$};

    % 3. Nhảy tần số F0: c_2 -> c_3 (chi phí = 0,35*|log2(f1/f2)|)
    \draw[wire_trans, draw=blue!80!black] (c_p_v1.east) -- (c_c_v2.west);
    \node[font=\scriptsize, fill=white, inner sep=2pt, rounded corners=2pt, draw=blue!50, text=blue!80!black, align=center] (lbl_nhay) at (4.0, 1.05) {Nhảy $F_0$:\\$0{,}35|\log_2\frac{f_1}{f_2}|$};
    \draw[draw=blue!60, thin, ->] (lbl_nhay.north) -- (4.0, 2.45);

    % Chuyển tiếp giữa khung t và t+1
    \draw[wire_trans] (c_c_uv.east) -- (c_n_uv.west);
    \draw[wire_trans] (c_c_v1.east) -- (c_n_v2.west);
    \draw[wire_trans] (c_c_v2.east) -- (c_n_v1.west);

    % Đường đi tối ưu Viterbi toàn cục (màu đỏ đậm có mũi tên)
    \draw[wire_viterbi] (c_p_v1.east) -- (c_c_v1.west);
    \draw[wire_viterbi] (c_c_v1.east) -- (c_n_v1.west);

    % -------------------------------------------------------------------------
    % ĐƯỜNG TRUY VẾT NGƯỢC (BACKTRACKING)
    % Mũi tên trỏ từ khung cuối N ngược về khung đầu (phải sang trái)
    % -------------------------------------------------------------------------
    \draw[draw=orange!90!black, ultra thick, dashed, ->] (14.0, 5.6) -- node[above=2pt, font=\scriptsize\bfseries, text=orange!90!black] {Truy vết ngược toàn bộ file ($N \to 1$)} (2.0, 5.6);

    % -------------------------------------------------------------------------
    % CHÚ THÍCH PHÂN TÁCH VITERBI VS MEDIAN-3
    % -------------------------------------------------------------------------
    \node[draw=red!70!black, fill=red!3, rounded corners=2pt, font=\scriptsize, align=left, inner sep=4pt] at (8.0, -1.6) {
        \textbf{Praat Path Finder (R1):} Lưu trữ toàn bộ ma trận $\delta, \psi$ kích thước $N \times 15$; bước tính xuôi cộng dồn chi phí chuyển tiếp,\\
        sau đó truy vết ngược từ khung cuối $N$ về khung 1 $\implies$ Bùng nổ RAM và độ trễ toàn file trên FPGA.
    };

\end{tikzpicture}%
}
